\section{Approach}
\label{sec:approach}

Figure~\ref{fig:overview} shows the steps of \tool. The main idea is to keep
every generated query as a typed derivation tree until it has been run, so
that oracles can rewrite single nodes of known type and position instead of
SQL text. The tree has four jobs: generation builds it, mutation re-derives
one subtree, the oracles rewrite one typed node, and triage reduces it by
replacing subtrees with their smallest derivation. The fuzzing loop (Steps
1--4) runs thousands of test cases per minute; the LLM loop (Step 5) runs a few
times per commit, before the campaign.

\begin{figure*}[t]
  \centering
  %% Figure: overview of the tool. Steps left to right, loop dashed.
%% Grayscale-safe (black lines, black!6 fill), \scriptsize text, numbered
%% steps matching the "Step N" subsections.
\begin{tikzpicture}[
    x=1cm, y=1cm, font=\scriptsize,
    box/.style={draw, rounded corners=2pt, align=center, inner sep=2pt,
                minimum height=9mm, text width=15mm, fill=white},
    outbox/.style={draw, align=left, inner sep=2pt, fill=black!6,
                text width=20mm, font=\tiny},
    flow/.style={-{Stealth[length=1.6mm]}, line width=0.5pt},
    back/.style={-{Stealth[length=1.6mm]}, line width=0.5pt, dash pattern=on 2pt off 1.5pt},
  ]
  \node[box, text width=10mm] (in) at (0,0) {input};
  \node[box] (s1) at (1.8,0) {\textbf{1} step one};
  \node[box] (s2) at (3.8,0) {\textbf{2} step two};
  \node[box] (s3) at (5.8,0) {\textbf{3} step three};
  \node[outbox] (o) at (3.8,-1.6) {output one\\output two};
  \draw[flow] (in) -- (s1);
  \draw[flow] (s1) -- (s2);
  \draw[flow] (s2) -- (s3);
  \draw[flow] (s3.south) |- (o.east);
  \draw[back] (s3.north) -- ++(0,0.35) -| node[pos=0.25, above] {repeat} (s2.north);
\end{tikzpicture}

  \caption{Overview of \tool. The top row is the fuzzing loop: test cases are
  derived from the typed grammar, oracles are attached as rewrites of the
  derivation tree, and the server runs each query under several settings.
  The bottom row runs before a campaign: an LLM turns a recent commit into
  grammar alternatives, and the \gate accepts them only when the server runs
  them and they reach the changed lines.}
  \label{fig:overview}
\end{figure*}

\subsection{Step 1: Derive Test Cases from a Typed Grammar}
\label{sec:grammar}

We wrote a grammar for the SQL that \pg~18 and master accept, in a small format
of our own (\pgg): 78 rules with 421 alternatives over 13 SQL types. A rule may
be parameterised by a type, as in \texttt{expr<T>}, so a derivation only
combines expressions that type-check. Weights set how often each alternative is
chosen, and \emph{builtins} written in Python handle what a context-free
grammar cannot: \texttt{@col(T)} picks a visible column of type $T$ in the
current scope, \texttt{@table} binds a fresh alias, and \texttt{@group\_keys}
makes the current query level grouped, after which only grouped columns and
aggregates are visible. The schema has 9 relations with partitions, inheritance,
a view, indexes of four access methods and NULLs in most columns.

A test case has up to two setup statements (DDL, DML, \texttt{ANALYZE},
\texttt{SET LOCAL} of a planner setting) followed by one or two queries. While
deriving, the engine wraps the tokens of every node in invisible markers, so
after rendering it knows the exact token span of each node. Half of the test
cases are new derivations; the others replay a corpus test case and re-derive one
subtree. A test case enters the corpus when it produces a new bug signature, a
new SQL feature, a new pair of features, or a new plan shape, the digest of the
node types in \texttt{EXPLAIN} output~\cite{ba2023qpg}.

\paragraph{Feedback and data.} Every alternative carries a multiplicative
weight that grows when a test case that used it enters the corpus and shrinks
when its statement fails (rate $\eta=0.1$, range $[0.2,5]$). Some literals come
from a \emph{data profile} (minimum, maximum and quantiles of every column,
type edges such as $-0$, NaN, $\pm$Infinity and the empty array), so predicates
select a real share of rows. Grammar-based generation is not new; the typing,
scoping builtins and node spans are what the later steps need.

\subsection{Step 2: Attach Oracles as Tree Rewrites}
\label{sec:oracles}

For each query, \tool picks up to one metamorphic check and stores it with
the test case as a pair of queries that must return the same multiset. All
checks are built on the derivation tree.

\paragraph{Partitioning and counting.} \emph{TLP on \texttt{WHERE}} picks a
\texttt{SELECT} whose result is used as a multiset (the query itself, a
derived table, a CTE body, a set-operation arm, or an \texttt{IN},
\texttt{EXISTS} or \texttt{LATERAL} subquery) and partitions its
\texttt{WHERE p} into $p$, \texttt{NOT p} and \texttt{p IS NULL}, combined
with \texttt{UNION ALL}. Because the rewrite edits one node, it also works in
a nested query, where text-based TLP does not apply. \emph{TLP on
\texttt{HAVING}} does the same; the unpartitioned side keeps \texttt{HAVING
TRUE}, because without \texttt{GROUP BY} a \texttt{HAVING} clause is what makes a
level a single-group aggregate. \emph{TLP on aggregates} equates a global
\texttt{MIN}, \texttt{MAX}, \texttt{SUM}, \texttt{COUNT}, \texttt{BIT\_*} or
\texttt{BOOL\_*} with the combination over the three partitions (the
\texttt{MIN} of the \texttt{MIN}s, the \texttt{SUM} of the \texttt{COUNT}s).
\emph{NoREC} compares \texttt{count(*)} with \texttt{WHERE p} against the sum
of \texttt{CASE WHEN p THEN 1 ELSE 0 END}.

\paragraph{Equivalent forms of one node.} \emph{Expression rewrites} replace
one typed expression $e$ by \texttt{COALESCE(e,e)}, \texttt{(SELECT e)},
\texttt{CASE WHEN TRUE THEN e END}, $e+0$ or $e \cdot 1$ (numbers),
\texttt{e || ''} (text), or \texttt{NOT NOT p}, \texttt{p AND TRUE},
\texttt{p OR FALSE} (predicates). \emph{Predicate rewrites} turn
\texttt{x IN (a, b)} into \texttt{(x = a) OR (x = b)}, which \pg~18 turns back
into \texttt{= ANY}, \texttt{a < b} into \texttt{NOT (a >= b)}, which
three-valued logic keeps exact, and \texttt{BETWEEN} into two comparisons.
\emph{NOT NULL conjunct} turns \texttt{WHERE p} into \texttt{WHERE p AND c IS
NOT NULL} for a NOT NULL column $c$ that $p$ uses, which targets the planner's
NOT NULL simplifications; it is skipped under outer joins and for generated
columns, where such a column can still be NULL. \emph{ORDER BY removal} drops
the statement's own \texttt{ORDER BY}, \emph{CTE materialisation} toggles
\texttt{MATERIALIZED}, and the \emph{table fence} turns one table reference
into \texttt{(SELECT * FROM t OFFSET 0)}, which stops the planner from pushing
predicates into it, removing it, or pruning its partitions.

\paragraph{Security and estimates.} For \emph{row-level security}, five tables
carry a policy defined in one file used by both the server and the rewrite.
The query runs as a restricted role, under the policies, and as a superuser
with every protected table $t$ replaced by \texttt{(SELECT * FROM t WHERE
\textit{policy})}; a difference means that the policies let a row through or
hide one they should not. Following CERT~\cite{ba2024cert}, adding the
\texttt{WHERE} clause to a query must not raise the planner's row estimate;
this check compares \texttt{EXPLAIN} estimates, and a violation is a
\emph{performance} candidate, reported apart from bugs.

A check is attached only if its rewrite is sound for that query
(Section~\ref{sec:problem}).

\subsection{Step 3: Run the Test Case Several Ways}
\label{sec:run}

The server runs each test case in a transaction with one savepoint per
statement and rolls it back at the end, so test cases cannot affect each other.
For each deterministic query \tool also runs (i)~both sides of the attached
check; (ii)~the query under two of 18 \emph{plan variants} (for example hash
join only, no hash aggregation, forced parallel plans, partitionwise join, no
eager aggregation, JIT), chosen at random per query; and (iii)~the whole test
case on a server of the last release, \pg~18.6, in a database of its own for
each worker (the \emph{cross-version oracle}). Results are compared as
order-insensitive digests that drop choices no plan controls: trailing zeros of
equal \texttt{numeric} values, the sign of float zero, the last digits of finite
floats in aggregates, and the text of equal \texttt{interval} values.

\subsection{Step 4: Classify, De-duplicate, Triage}
\label{sec:triage}

The oracle reads the statement results and the server log. A \texttt{TRAP} line
is an assertion failure; a lost connection without one is a crash; an error of
class \texttt{XX} is an internal error; a confirmed mismatch of any digest is a
wrong result (a metamorphic mismatch counts only if it repeats when both sides
run again). Crash evidence counts only if it carries the process id of the test
case's own backend or its parallel workers. Each candidate gets a
\emph{signature}: the assertion location, the normalised backtrace or error
message, or, for wrong results, the oracle and the shape of the query; only the
first test case per signature is saved.

Triage reproduces the candidate three times on a fresh connection and reduces
it in two stages. \emph{Tree reduction} works on the derivation trees that every
failure of the grammar arm stores: it drops a statement, drops one repetition,
or replaces a subtree by the smallest derivation of the same rule and type, and
replays the trees through the generator, which re-binds names, aliases and
types. Every candidate is therefore valid SQL, and the server's time goes to the
question whether the failure remains. \emph{Token reduction} follows: delta
debugging~\cite{zeller2002dd} removes statements, then clauses, list items and
parenthesised groups level by level over the parenthesis
tree~\cite{misherghi2006hdd}, then single tokens. A step is kept only if the same
failure occurs (same assertion, normalised message, or oracle and plan variant);
for a metamorphic candidate both queries are reduced together, so they stay a
valid pair.

\paragraph{Localization.} For a reduced wrong result, \tool runs the queries
under every plan variant. If a variant removes the disagreement, the bug is in
the code path that setting switches, and the report names the setting and its
source files; if none does, the cause is an always-applied rewrite, constant
folding, a planner support function, or expression evaluation. The output is a
short reproducer, the rows each side returned, the localization and a draft
report in the shape of the \pg bug form.

\subsection{Step 5: Extend the Grammar from Recent Commits}
\label{sec:llm}

\paragraph{Step 5a: propose.} \tool lists the commits of \pg master that change
the optimiser, executor, parser or SQL-callable functions and add regression
tests. For each it builds a digest: the message, the changed functions, the SQL
the commit added to the tests, and the \emph{diff-coverage targets}, the changed
lines of each C file. An LLM gets the digest with a fixed, cached system prompt
(the \pgg format, the builtins, the conventions of the grammar, the schema and
the whole current grammar) and answers with structured output: a rationale, a
fragment of new alternatives, one to five seed test cases, the features it
targets, and up to three \emph{equivalence rewrites}, templates over one
expression of a given type such as \texttt{CASE WHEN \{e\} IS NULL THEN NULL
ELSE \{e\} END} that the model claims equal \texttt{\{e\}} for every value. The
commit miner never shows the model a fix that the benchmark of
Section~\ref{sec:rq3} reverts.

\paragraph{Step 5b: validate.} No model output reaches the campaign unless the
server agrees. The \gate checks in order: (i)~\emph{parse}: the fragment parses
with the grammar, adds no types and uses only known builtins;
(ii)~\emph{reachability}: in a sample of 150 test cases with the new alternatives
boosted, they are used at least five times; (iii)~\emph{validity}: at least 60\%
of the statements that use them run without error; (iv)~\emph{seeds}: each seed
runs; and (v)~\emph{diff coverage}: on a gcov build, the sample executes more of
the commit's changed lines than the same sample without the fragment. A failed
check sends its findings (the parse error, the most frequent error messages, the
uncovered lines) back to the model for up to three repair rounds; every round is
a ledger line with its tokens, cost and the gate's report.

A proposed rewrite is a claim about semantics, so the gate tests it twice. The
server evaluates \texttt{rewrite(v) IS NOT DISTINCT FROM v} for NULL, the edge
values of the type and quantiles of the data; then \tool applies the rewrite
alone to the expressions of 150 test cases, with confirmation re-runs. The
rewrite is accepted only if no value and no query shows a difference, and it
then becomes an expression rewrite of Step~2. Neither test varies the
collation, so the gate also rejects a template with an explicit
\texttt{COLLATE} (Section~\ref{sec:rq4}).
